%% ================================================================================================
%%  8. Discussion
%%  Notes: research-agenda.md §8 (open questions), author-tasks.md
%% ================================================================================================
\section{Discussion}\label{sec:discussion}

\subsection{How firm are the conclusions?}

Three conclusions rest on theorems and are not in doubt: no procedure is immune to manipulation; the
majority principle and a minority veto conflict in some profiles; and individual vetoes with the status
quo as default deadlock a large assembly. Two rest on a broad and consistent literature: Condorcet methods
achieve the best welfare among common rules when preferences have structure, and an open agenda calls for
independence of clones. The recommendation of a blocking quota of one third is a judgement on a
trade-off, supported by simulations on stylised electorates; the shape of the trade-off is robust across
our cultures, the precise numbers are not.

\subsection{Limits of this note}

The note considers a single decision. Repeated decisions raise questions it does not treat: today's
outcome becomes tomorrow's status quo \citep{DziudaLoeper2016}, vetoes can be traded or saved
\citep{Casella2005}, and reputations form. It takes preferences and proposals as given, whereas
deliberation changes both, and the IETF's practice suggests that much of the value of a consensus process
lies in how objections are discussed rather than counted \citep{Resnick2014rfc7282}. All voters have equal
weight; weighted assemblies (shareholders, federations) would need weighted counts and a weighted quota,
which the procedure accommodates but we have not analysed. Finally, the literature search behind the
Consensus review, and our own additions, focus on economics, political science and computational social
choice; work in management science on group decision support, and in law on supermajority requirements,
is not covered.

\subsection{Open questions}

\begin{itemize}
  \item \textbf{A clone-proof proportional veto?} The proportional veto core protects minorities in a
  way that no coalition can undo, but fails independence of clones; the blocking quota is clone-proof
  but open to strategic blocking (Propositions~\ref{prop:pvc-clones} and~\ref{prop:blocking}). Whether
  some rule combines independence of clones, veto power proportional to size and a non-empty core, or
  whether the three are incompatible, is open as far as we know.
  \item \textbf{Fast procedures with hard vetoes.} One- and two-round proposal procedures reach the
  Condorcet winner quickly under single-peaked preferences \citep{Gersbach2025propose,GersbachValkanova2024};
  their behaviour with binding vetoes and several dimensions is unknown (a gap also noted by the Consensus
  review).
  \item \textbf{Strategic comparison of compromise rules.} Fallback bargaining, shortlisting, veto-rank and
  proportional veto have been analysed under different assumptions; a common framework with strategic
  voters and many parties is missing \citep{BarberaCoelho2022,CaillouxNapolitanoSanver2022}.
  \item \textbf{Behaviour.} Do named and reasoned vetoes reduce strategic blocking in practice? The
  evidence on responsibility in committees \citep{BreitmoserValasek2024} and on tacit consent
  \citep{Krick2017} suggests so; a laboratory test of \QVC\ with and without such safeguards would settle
  it.
  \item \textbf{Real ballots.} The conclusions should be checked on observed ranked ballots with
  acceptability lines, of which Debian's general resolutions are a public source
  \citep{Debian2024constitution,MatteiWalsh2013}.
\end{itemize}

\subsection{Putting the procedure to work}

For an organisation that wants to adopt it, the procedure fits on one page.
\begin{enumerate}
  \item \textbf{Fix the parameters in advance}: the blocking quota (one third by default; as a number of
  voters, $\lceil n/3 \rceil$), whether the status quo is protected (by default not), the length of the
  nomination period, and the rule for the optional revision round.
  \item \textbf{Nomination}: anyone may propose any number of options; the status quo, or ``none of the
  above'', is on the ballot. Identical proposals may be merged, but variants need not be.
  \item \textbf{Ballot}: rank or grade every option, ties allowed, and mark the unacceptable ones; a veto
  is named and gives a reason.
  \item \textbf{Count} with Algorithm~\ref{alg:qvc}; publish the veto counts, the admissible options, the
  pairwise matrix and the winner. The count can be checked by any member with the reference code
  (\texttt{simulations/python/voting.py}).
  \item \textbf{Revision round}, only if every option was vetoed by the quota: reopen nominations briefly,
  vote again, and accept the second result as final.
\end{enumerate}
